% !TEX root = print.tex

% In this file you should put the actual content of the blueprint.
% It will be used both by the web and the print version.
% It should *not* include the \begin{document}
%
% If you want to split the blueprint content into several files then
% the current file can be a simple sequence of \input. Otherwise It
% can start with a \section or \chapter for instance.

\begin{abstract}
    Our main objective is the verification of an implementation of Asynchronous Byzantine
    Agreement (ABA). The proof of correctness consists of building a simulation relation between
    the implementation and a transition system encoding its specification. The simulation
    ensures that the encoding of the specification implies the correctness of the
    original one. However, this implementation is complex; thus, building the simulation in one
    go is intractable. Fortunately, the implementation is designed as the interlocking of
    multiple sub-protocols (GBCA, WCC, RSD, Gather, AVSS, BRB). This architecture allows us to
    decompose the construction of the overall simulation into smaller steps. Each implementation
    of sub-protocols will be simulated by the transition system of its specification. Thus, when 
    building the simulation for a protocol using sub-protocols, the sub-implementations will be 
    replaced by their specification, making the construction easier.
\end{abstract}

\section{Preliminaries}

\subsection{Segala's Framework}

\begin{definition}[Probabilistic Labelled Transition System]
    \lean{PLTS}
    %\leanok

    A probabilistic labelled transition system (PLTS) is a tuple $\pcal := (Q,q_\star,L,T)$
    where $Q$ is a (countable) set of states,
    $q_\star \in Q$ is the initial state,
    $L$ is a (countable) set of labels and
    $T \subseteq Q \times (L \cup \{\tau\}) \times \dcal(Q)$ is a set of probabilistic transitions.
\end{definition}

A partial execution of $\pcal$ is a sequence
$q_0.l_0.q_1.l_1\cdots \in (Q.L)^*.Q \cup (Q.L)^\omega$ such that, for every $n \in \nat$, there
exists $\mu_n \in \dcal(Q)$ with $(q_n,l_n,\mu_n) \in T$ and $q_{n+1} \in \bar{\mu_n}$. The set of
finite executions is denoted $\ecal^*$. A partial execution is called an execution if
$q_0 = q_\star$. The trace of a partial execution is its projection onto $L$. The set of traces of
$\pcal$ is denoted $\tcal$. A scheduler $s$ is a function from $\ecal^*$ to $\dcal(T \cup \{\bot\})$
such that, for all $e \in \ecal^*$ and $t \in \bar{s(e)}$, $t = \bot$ or $src(t) = last(e)$. The
set of schedulers is denoted $Sch$.

\begin{definition}[Partial Probabilistic Execution]
    \lean{Partial Probabilistic Execution}
    %\leanok

    Given $\mu \in \dcal(Q)$ and $s \in Sch$, one can define the partial probabilistic execution
    induced by $(\mu,s)$ as the function $\efrak_{\mu,s}$ from $\ecal^*$ to $[0,1]$ such that if
    $e \in Q$ then $\efrak_{\mu,s}(e) = \mu(e)$; otherwise, $e = e'.l.q$ and
    $\efrak_{\mu,s}(e) =
    \efrak_{\mu,s}(e')\cdot\sum_{t \in T \mid lbl(t) = l}s(e')(t)\cdot out(t)(q)$.
\end{definition}

A partial probabilistic execution is called a probabilistic execution if it is induced by a pair
$(\mu,s)$ such that $\mu = \delta_{q_\star}$. If $s \in Sch$ terminates almost surely from $\mu$
(i.e. $\sum_{e \in \ecal^*} \efrak_{\mu,s}(e)\cdot s(e)(\bot) = 1$) then
$\qfrak_{\mu,s} = last^{-1}\circ\efrak_{\mu,s}$ is a well-defined probability distribution over
states. We denote $\mu \xRightarrow{l} \mu'$ if and only if there exists $s \in Sch$ such that
$\delta_l = \tfrak_{\mu,s}$ and $\mu' = \qfrak_{\mu,s}$. Partial probabilistic executions induce a
measure on the $\sigma$-field induced by cones of partial executions, and the trace function is
measurable from the $\sigma$-field induced by cones of partial executions to the one induced by
cones of traces. Thus, $\tfrak_{\mu,s} = tr^{-1}\circ\efrak_{\mu,s}$ is well-defined and called a
trace distribution. The set of trace distributions achievable by $\pcal$ is defined as those
induced by a probabilistic execution of $\pcal$ and denoted $ \text{tdist}(\pcal)$.
A hyper-transition of a PLTS $\pcal$ is a triple $(\mu,l,\mu')$ such that, for all $q \in \bar\mu$,
there exists $(q,l,\mu_q) \in T$ such that $\mu' = \sum_{q \in Q}\mu(q)\cdot\mu_q$.
The set of hyper-transitions is denoted $H$.
The trace distribution pre-order $\pcal \subseteq^t \pcal'$ is defined by
$ \text{tdist}(\pcal) \subseteq  \text{tdist}(\pcal')$.

\begin{definition}[Parallel Composition of PLTS]
    \lean{Parallel Composition}
    %\leanok

    Given two PLTSs $\pcal$ and $\pcal'$ and $S \subseteq L \cap L'$, the composed PLTS
    $\pcal \parallel_S \pcal'$ is defined as
    $(Q \times Q', (q_\star, q_\star'), L \cup L', T \parallel_S T')$ where
    $((q,q'), l, \mu'') \in T \parallel_S T'$ if and only if
    \begin{itemize}
        \item $l \in S$ and there exist $(q,l,\mu) \in T$ and $(q',l,\mu') \in T'$ such that
        $\mu \otimes \mu' = \mu''$; or
        \item $l \not\in S$ and there exists $(q,l,\mu) \in T$ such that
        $\mu \otimes \delta_{q'} = \mu''$; or
        \item $l \not\in S$ and there exists $(q',l,\mu') \in T'$ such that
        $\delta_q \otimes \mu' = \mu''$.
    \end{itemize}
\end{definition}

The trace distribution pre-order is not stable under $\parallel_S$: $\pcal \subseteq^t \pcal'$ does
not imply $\pcal \parallel_S \pcal'' \subseteq^t \pcal' \parallel_S \pcal''$
\cite{Segala1995Compositional}.
Therefore, the largest pre-order included in $\subseteq^t$ stable under $\parallel_S$ is considered:
$\pcal \sqsubseteq^t \pcal'$ if and only if, for all PLTSs $\pcal''$ and
$S \subseteq L \cap L' \cap L''$,
$\pcal \parallel_S \pcal'' \subseteq^t \pcal' \parallel_S \pcal''$.

\begin{lemma}
    \label{lem:precong}
    If $\pcal \sqsubseteq^t \pcal'$ then, for all PLTSs $\pcal''$ and
    $S \subseteq L \cap L' \cap L''$,
    $\pcal \parallel_S \pcal'' \sqsubseteq^t \pcal' \parallel_S \pcal''$
\end{lemma}
\begin{proof}
    This is a direct consequence of the associativity of $\parallel_S$.
\end{proof}

\begin{definition}[Abstraction Operator]
    Given a PLTS $\pcal$ and $I \subseteq L$, the abstract PLTS $\tau_I(\pcal)$ is obtained
    by replacing all labels in $I$ by $\tau$.
\end{definition}

\begin{lemma}
    \label{lem:abstract}
    If $\pcal \sqsubseteq^t \pcal'$ then, for all $I \subseteq L$,
    $\tau_I(\pcal) \sqsubseteq^t \tau_I(\pcal')$.
\end{lemma}

Given $R \subseteq A \times \dcal(B)$, the probabilistic lifting of $R$, denoted $\dcal(R)$, is
defined by $(\mu_A,\mu_B) \in \dcal(R)$ if and only if there exists a witness $w : R \rightarrow 
[0,1]$ such that $\mu_A = a \mapsto \sum_{\mu \in \dcal(B)}w(a,\mu)$ and $\mu_B = \sum_{(a, \mu) 
\in R}w(a,\mu)\cdot\mu$.

\begin{definition}[Probabilistic Forward Simulation]
    \lean{Probabilistic Forward Simulation}
    %\leanok

    Given two PLTSs $\pcal$ and $\pcal'$, a probabilistic forward simulation of $\pcal$ by $\pcal'$
    is a relation $\scal \subseteq Q \times \dcal(Q')$ such that
    $(q_\star,\delta_{q_\star'}) \in \scal$ and, for all $(q,\mu') \in \scal$ and $(q,l,\mu) \in T$,
    there exists $\mu' \xRightarrow{l} \mu''$ such that $(\mu,\mu'') \in \dcal(\scal)$.
\end{definition}

\begin{theorem}[from \cite{Segala1995Compositional}]
    \label{thm:Segala}
    Given two PLTSs $\pcal$ and $\pcal'$, $\pcal \sqsubseteq^t \pcal'$ if and only if there exists a
    probabilistic forward simulation of $\pcal$ by $\pcal'$.
\end{theorem}
\begin{proof}
    The proof is non-trivial and is detailed in Segala's PhD thesis.
    Our approach to this proof is detailed in Appendix \ref{app:ProofSegala}.
\end{proof}

\subsection{Non-probabilistic Systems}

Non-probabilistic systems can be considered as special instances of PLTSs in which all transitions
lead to a Dirac distribution. By identifying each Dirac distribution with the unique state in its
support, we may assume that non-probabilistic systems are modelled by LTSs, which leads to several
simplifications. In particular, probabilistic forward simulation is equivalent to a simpler
simulation relation. On an LTS, we write $q \xRightarrow{l} q'$ if there exists a partial execution
from $q$ to $q'$ whose trace is $l$.

\begin{definition}[Weak Simulation]
    \lean{Weak Simulation}
    %\leanok

    Given two LTSs $\pcal$ and $\pcal'$, a weak simulation of $\pcal$ by $\pcal'$ is a
    relation $\scal \subseteq Q \times Q'$ such that $(q_\star,q_\star') \in \scal$ and,
    for all $(q,q') \in \scal$ and $(q,l,p) \in T$,
    there exists $q' \xRightarrow{l} p'$ such that $(p,p') \in \scal$.
\end{definition}

\begin{lemma}
    \label{lem:weaksim}
    Given two LTSs $\pcal$ and $\pcal'$, $\pcal \sqsubseteq^t \pcal'$ if and only if there exists a
    weak simulation of $\pcal$ by $\pcal'$.
\end{lemma}
\begin{proof}
    According to Theorem \ref{thm:Segala}, it suffices to prove that
    there exists a weak simulation of $\pcal$ by $\pcal'$ if and only if
    there exists a probabilistic forward simulation of $\pcal$ by $\pcal'$. For one direction,
    it suffices to note that any weak simulation between LTSs is a probabilistic forward simulation.
    For the other direction, if $\scal$ is a probabilistic forward simulation of $\pcal$ by $\pcal'$
    then $\{(q,q') \mid \exists (q,\mu) \in \scal, q' \in \bar\mu\}$ is a weak simulation of $\pcal$
    by $\pcal'$.
\end{proof}

There is a natural injection of LTSs into PLTSs, but also a natural projection of PLTSs onto LTSs.
It consists in transforming all probabilistic non-determinism into classical non-determinism.
The projection of a PLTS $\pcal$, denoted $\rho(\pcal)$, is the LTS whose transitions are defined
by $\{(q,l,q') \mid \exists \mu \in \dcal(Q), (q,l,\mu) \in T \wedge q' \in \bar\mu\}$.
$\rho$ is trivially the identity on LTSs, and $\rho$ preserves the set of executions and hence the
set of traces. As a result, the notion of weak simulation can be lifted to PLTSs: $\scal$ is a weak
simulation of $\pcal$ by $\pcal'$ if and only if it is a weak simulation of $\rho(\pcal)$ by
$\rho(\pcal')$. Since $\rho$ is a projection onto LTSs, this is consistent with the previous
definition.

\subsection{Preserving Liveness}

Liveness properties should not be tested with respect to the set of trace distributions and the set 
of traces of a PLTS because these sets are downward-closed and, in particular, they respectively 
always contain the Dirac distribution of the empty trace and the empty trace. Therefore, against 
such sets, liveness properties like termination are not satisfied by any PLTS. This may seem absurd
but it hints at an implicit progress assumption made when evaluating liveness properties. This
progress assumption could be stated as: if a transition is enabled then a transition must fire.
Thus, it discards all executions that are finite and whose last state is not a deadlock. Executions 
satisfying the progress assumptions are said to be maximal and traces of maximal executions are
also called maximal. Non-probabilistic liveness properties should be tested with respect to the set of
maximal traces. Similarly, a scheduler is said to be maximal if it schedules $\bot$ with positive 
probability only from maximal executions and a trace distribution is said maximal if it is induced
by a maximal scheduler. Probabilistic liveness properties should be tested against maximal trace 
distribution.

In the setting of distributed systems suffering from Byzantine failures, another assumption is
necessary in order to define a sensible evaluation of liveness properties. Indeed, since Byzantine
processes behave arbitrarily, they can monopolise the execution preventing correct processes from
taking a step. As a result, a fairness assumption is introduced stating that, in any infinite 
execution, correct processes must be scheduled infinitely often. To model it formally, a fairness 
function $\text{fair} : T \rightarrow \bool$ is introduced which indicates which transitions are 
fair (e.g. transition of correct processes) and which are unfair (e.g. transitions of Byzantine 
processes). However, Byzantine failures also interfere with the notion of deadlock as they are 
allowed to crash at any time, thus, executions may stop if no fair transition is enabled. We call a
fair deadlock any state whose set of enabled transitions is included in the set of unfair 
transition. An execution is said to be fair if (a) it is finite and ends in a fair deadlock, or (b)
it contains infinitely many fair transitions. A trace is said fair if it is induced by a fair 
execution. Non-probabilistic liveness properties will be evaluated with respect to the set of fair 
traces. Similarly, a scheduler is fair if (a) it schedules $\bot$ with positive probability only 
from fair deadlocks and (b) all infinite executions consistent with it are fair. A trace 
distribution is fair if it is induced by a fair scheduler. Probabilistic liveness properties will 
be evaluated against the set of fair trace distributions.

The set of fair traces is denoted $\tcal^f$ and the set of fair trace distributions is denoted 
$\text{tdist}^f(\pcal)$. Unfortunately, inclusion of $\tcal^f$ is not preserved by parallel
composition because synchronisation issues could create new deadlocks. However, in our case, it 
never happens because synchronisation will be used to call sub-processes locally. The modelling 
solution is to use the input-output approach of \cite{InputOutput} where the set of 
labels is partitioned between input labels $In$ and output labels $Out$. We require that the input 
labels are always enabled: $\forall q \in Q, l \in In, \exists \mu \in \dcal(Q), (q,l,\mu) \in T$. 
This requirement is easily fulfilled by adding self-looping transition labelled by each input label
to all states. Moreover, we restrict composition to PLTSs whose outputs labels are distinct. This 
restriction is not problematic either because labels will encode the protocol APIs. Under such 
assumptions, parallel composition does not create new deadlocks and so the inclusion of $\tcal^f$ 
is preserved by composition. As for $\subseteq^t$, inclusion of $\text{tdist}$ is not preserved by 
parallel composition, and the problem cannot be solved by harmless restrictions. Thus, the fair 
trace distribution pre-congruence is defined as: 
\[
    \pcal \sqsubseteq^f \pcal' \iff \forall \pcal'', S \subseteq L, \text{tdist}^f(\pcal \parallel_S \pcal'') \subseteq  \text{tdist}^f(\pcal' \parallel_S \pcal'') 
\]
It enjoys the same algebraic properties as $\sqsubseteq^t$.

Probabilistic forward simulation and weak simulation preserve $\sqsubseteq^t$, but not 
$\sqsubseteq^f$. Looking at the folklore construction made in Lemma \ref{lem:weaksim}, a fair 
execution can be simulated by an unfair execution. If the execution ends with a fair deadlock then
the simulation does not guarantee anything about the last state of the simulating execution.
Moreover, internal fair transitions can be elided, thus, an infinite fair execution can be simulated
by an execution with finitely many fair transitions. We call an execution a fair divergence if it is
fair and its trace is $\tau$. The diagnostic is the following: both simulations do not simulate
fair divergences properly. We will say that a weak simulation $\scal$ is fair if, for all $(q,p) 
\in \scal$, if $q$ fairly diverges then $p$ fairly diverges. However, this simple solution has a
major problem: proving that a state does not fairly diverge is hard. Fortunately, 
\cite{Dongol2022Weak} proposes a solution to a similar problem: equip the simulation with a ranking
function. We adapt their contribution to our setting. 

\begin{definition}[Fair Weak Simulation]
    \lean{FairWeakSimulation}
    %\leanok

    Given two LTSs $\pcal$ and $\pcal'$, a fair weak simulation of $\pcal$ by $\pcal'$ is a 
    relation $\scal \subseteq Q \times Q'$ equipped with a well-founded preorder $\le$ on $Q$ such 
    that $(q_\star,q_\star') \in \scal$ and, for all $(q,q') \in \scal$ and $(q,l,p) \in T$, there 
    exists $q' \xRightarrow{l} p'$ such that $(p,p') \in \scal$. Furthermore, if $l = \tau$ and 
    $q' \xRightarrow{\tau} p'$ contain no fair transition then if $(q,l,p)$ is fair then $q > p$; 
    else $q \ge p$. Moreover, for all $(q,q') \in \scal$, if $q$ is a fair deadlock then $q'$ 
    fairly diverges.
\end{definition}

For probabilistic forward simulation, the problem is more complicated because the construction of
the simulating scheduler is not so clear (see Appendix \ref{app:ProofSegala}). However, the core
idea stays the same even though it is more convoluted.

\begin{definition}[Fair Probabilistic Forward Simulation]
    \lean{FairProbForwardSimulation}
    %\leanok

    Given two LTSs $\pcal$ and $\pcal'$, a fair probabilistic forward simulation of $\pcal$ by 
    $\pcal'$ is a relation $\scal \subseteq Q \times \dcal(Q')$ equipped with a well-founded 
    preorder $\le$ on $Q$ such that $(q_\star,\delta_{q_\star'}) \in \scal$ and, for all $(q,\mu') 
    \in \scal$ and $(q,l,\mu) \in T$, there exists $\mu' \xRightarrow{l} \mu''$ such that 
    $(\mu',\mu'') \in \dcal(\scal)$. Furthermore, let $s$ be the scheduler inducing $\mu' 
    \xRightarrow{\tau} \mu''$ and $w$ the witness of $(\mu',\mu'') \in \dcal(\scal)$, if $l = \tau$
    and there exist $q' \in \bar\mu' \xRightarrow{\tau} p'$ consistent with $s$ which contain no 
    fair transition then, for all $p \in \bar\mu$, if there exists $\mu_p \in \dcal(Q')$ such that
    $w(p,\mu_p) > 0$ and $\mu_p(p') > 0$ then if $(q,l,\mu)$ is fair then $q > p$; else $q \ge p$.
    Moreover, for all $(q,\mu') \in \scal$, if $q$ is a fair deadlock then there exists a fair 
    scheduler $s$ such that $\tfrak_{\mu',s} = \delta_\tau$.
\end{definition}

\begin{theorem}
    \label{thm:Segalafair}
    Given two PLTSs $\pcal$ and $\pcal'$, if there exists a fair probabilistic forward simulation 
    of $\pcal$ by $\pcal'$ then $\pcal \sqsubseteq^f \pcal'$.
\end{theorem}
\begin{proof}
    [TODO]
\end{proof}

\begin{lemma}
    \label{lem:fairweaksim}
    Given two LTSs $\pcal$ and $\pcal'$, if there exists a fair weak simulation of $\pcal$ by 
    $\pcal'$ then $\pcal \sqsubseteq^f \pcal'$.
\end{lemma}
\begin{proof}
    It suffices to note that a fair weak simulation between LTSs is a fair probabilistic forward
    simulation.
\end{proof}

\subsection{Partial Information}

We encode partial information by restrictions on the set of schedulers.
Intuitively, if the scheduler does not have access to all information then it might
not be able to distinguish some states or some labels; thus, some distinct
executions might look the same to it and it should schedule the same
probability distribution of transitions.

\begin{definition}[Adversary]
    An adversary is a pair of observation functions $(\ocal_Q, \ocal_L)$ which
    respectively map states and labels to state observation and label observation
    such that, for all $t,t' \in T$, if $src(t) = src(t')$ and
    $\ocal_L(lbl(t)) = \ocal_L(lbl(t'))$ then $t = t'$.
\end{definition}

The injectivity of $\ocal_L$ on the set of transitions enabled in a given
state ensures that partial information does not prevent the adversary from resolving
non-determinism. An adversary $(\ocal_Q,\ocal_L)$ induces an observation $\ocal_{\ecal}$
on executions defined by
$\ocal_{\ecal}(q_0.l_0\cdots) = \ocal_Q(q_0).\ocal_L(l_0)\cdots$.
A scheduler $s$ is consistent with an adversary if, for all $e,e' \in \ecal^*$,
$\ocal_{\ecal}(e) = \ocal_{\ecal}(e') \implies s(e) = s(e')$. The omniscient adversary
is encoded by setting the observation functions to the identity.

Our simulations are sound for partial-information adversaries because they were
designed to be sound for omniscient adversaries. However, there is a completeness
gap as partial information makes more reductions sound. Our chosen solution is to
build a PLTS on which the omniscient adversary is equivalent to the
partial-information adversary on the original PLTS. This is the belief construction.

\begin{definition}[Belief PLTS]
    Given a PLTS $\pcal$ and an adversary $\acal :=(\ocal_Q,\ocal_L)$, the belief PLTS
    $\pcal_{\acal}$ is defined as the tuple
    $(\{\mu \in \dcal(Q) \mid \ocal_Q(\bar\mu) = \{\cdot\} \},
    \delta_{q_\star}, \ocal_L(L), T_{\acal})$
    where
    \[
        \mu \xrightarrow{l} \mu' \in T_ {\acal} \iff
        \exists \mu \xrightarrow{l} \mu'' \in H,
        \mu' = (q \mapsto
        \frac{\mathbb{1}_{\ocal_Q(q) = o} \cdot \mu''(q)}{\mu''(\{q \mid \ocal_Q(q) = o\})})
        \mapsto \mu''(\{q \mid \ocal_Q(q) = o\})
    \]
\end{definition}

\begin{theorem}
    \label{thm:partialinfosound}
    Given a PLTS $\pcal$ and an adversary $\acal$, the set of trace distributions
    of $\pcal$ induced by a scheduler consistent with $\acal$ is equal to $ \text{tdist}(\pcal)$.
\end{theorem}

When considering the composition of two PLTSs $\pcal$ and $\pcal'$ scheduled by
their respective adversaries $\acal$ and $\acal'$, we need to define a composite
adversary. Given $\pcal$ and $\pcal'$ two PLTSs and $\acal$ and $\acal'$ two adversaries
(one for each PLTS) and $S \subseteq L \cap L'$, we say that $\acal$ and $\acal'$ are
compatible on $S$ if (i), for all $l \in S$, $\ocal_L(l) = \ocal'_L(l)$;
(ii) $\ocal_L(L \setminus S) \cap \ocal_L(S) = \emptyset$ and
(iii) $\ocal'_L(L' \setminus S) \cap \ocal'_L(S) = \emptyset$.
The three requirements state that the adversaries of each PLTS get the same
information from the synchronised labels and distinguish synchronised labels
from the others. If $\acal$ and $\acal'$ are compatible on $S$, we can define
$\acal \parallel_S \acal'$ as
\[
    ((q,q') \mapsto (\ocal_Q(q),\ocal'_Q(q')),
    l \mapsto
    \begin{cases}
        \ocal_L(l) \text{ if } l \in L \\
        \ocal_L'(l) \text{ if } l \in L'
    \end{cases})
\]

\begin{theorem}
    \label{thm:partialinfocompo}
    Given two PLTSs $\pcal$ and $\pcal'$ and their respective adversary $\acal$
    and $\acal'$ and $S \subseteq L\cap L'$, if $\acal$ and $\acal'$ are compatible
    on $S$ then $\pcal_{\acal} \parallel_S \pcal'_{\acal'}$ and
    $(\pcal \parallel_S \pcal')_{\acal \parallel_S \acal'}$ are isomorphic.
\end{theorem}
\begin{proof}
    For clarity, let us denote $\pcal_1 := \pcal_{\acal} \parallel_S \pcal'_{\acal'}$ and
    $\pcal_2 := (\pcal \parallel_S \pcal')_{\acal \parallel_S \acal'}$.
    First, note that $Q_1$ is in $\dcal(Q) \times \dcal(Q')$
    whereas $Q_2$ is in $\dcal(Q \times Q')$ and there is no bijection between
    the two. However, one can prove that $Q_2$ is actually included in
    $\{\mu \in \dcal(Q \times Q') \mid \mu := \mu' \otimes \mu''\}$, thus,
    $h : (\mu,\mu') \mapsto \mu \otimes \mu'$ is a bijection between $Q_1$ and $Q_2$.
    It remains to prove the morphism part which consists in unfolding definitions.
\end{proof}

\section{List of the Protocols}

\subsection{Main Protocol: Asynchronous Byzantine Agreement (ABA)} (taken from \cite{Abraham2022Efficient})
Processes start with a bit $\{0,1\}$ and may return a bit.
a program implements ABA if it satisfies three properties: Validity, Agreement, Almost-Sure Termination.
\begin{itemize}
    \item Validity: if a correct process returns $b$ then a correct process had $1-b$ as input.
    \item Agreement: if two correct processes return, then their output are equal.
    \item Termination: if $n-f$ correct processes take part to the protocol then, with probability 1, all correct processes will eventually return.
\end{itemize}

\begin{protocol}\label{prot:SpecABA}
    \lean{ABA.Spec}
    %\leanok
    \begin{algorithm}
        \caption{Ideal Implementation of ABA}
        \begin{algorithmic}
            \State \textbf{Input:} $b \in \{0,1\}$
            \medskip
            \State \textbf{Code for every process:}
            \State send $\langle input, b\rangle$ to the third party 
            \State \textbf{upon} receiving $\langle output, b\rangle$ from third party 
            \State \quad return $b$
            \medskip
            \State \textbf{Code for third party:}
            \State $bind,call \gets \bot,(\bot)^{Id}$
            \State \textbf{upon} receiving $\langle input, b\rangle$ from a process $id$ and $bind = \bot$:
            \State \quad $call(id) \gets b$
            \State \textbf{upon} $|F| + \{id \in Id \setminus F \mid call(id) = b\} \ge n-f$:
            \State \quad $bind,call \gets b,(\bot)^{Id}$
            \State \quad \textbf{if} $\{id \in Id \setminus F \mid call(id) = 1-b\} = 0$:
            \State \quad\quad send $\langle output, b\rangle$ to all
            \State \textbf{upon} $call(id) = \bot$ and $bind \ne \bot$:
            \State \quad $call \gets bind \textbf{ or } \mu$ \Comment{$\mu$ can be any non-Dirac probability distriution over $\{0,1\}$}
        \end{algorithmic}
    \end{algorithm}
\end{protocol}

\begin{protocol}\label{prot:ImplABA}
    \lean{ABA.Impl}
    %\leanok
    \begin{algorithm}
        \caption{Implementation of ABA taken from \cite{Abraham2022Efficient}}
        \begin{algorithmic}
            \State \textbf{Input:} $b \in \{0,1\}$
            \medskip
            \State $r \gets 1$
            \State $decided \gets \bot$
            \Loop
                \State $(b, g) \gets \textsc{GBCA}_r(b)$ \Comment{$b \in \{0,1,\bot\}$, $g \in \{A, B, C\}$}
                \State $c \gets \textsc{WCC}_r$ \Comment{common coin for round $r$}
                \If{$b = \bot$}
                    \State $b \gets c$
                \ElsIf{$g = A$}
                    \State $decided \gets b$ \Comment{ keep participating to help others}
                \EndIf
                \State $r \gets r + 1$
            \EndLoop
            \State \Comment{TODO: verify termination rule}
        \end{algorithmic}
    \end{algorithm}
\end{protocol}

\begin{definition}\label{def:ABAConcrete}
    \lean{ABA.Concrete}
    %\leanok
    \uses{prot:ImplABA, prot:ImplGBCA, def:WCCConcrete}
    $\mathsf{ABA.Concrete} := \mathsf{ABA.Impl}[\mathsf{GBCA.Impl}, \mathsf{WCC.Impl}]$, the instantiation of $\mathsf{ABA.Impl}$ with concrete implementations of the sub-protocols.
\end{definition}

\begin{definition}\label{def:ABAHybrid}
    \lean{ABA.Hybrid}
    %\leanok
    \uses{prot:ImplABA, prot:SpecGBCA, prot:SpecWCC}
    $\mathsf{ABA.Hybrid} := \mathsf{ABA.Impl}[\mathsf{GBCA.Spec}, \mathsf{WCC.Spec}]$, the instantiation of $\mathsf{ABA.Impl}$ with ideal specifications of the sub-protocols.
\end{definition}

\subsection{ABA Sub-Protocols}

\paragraph{Graded Binding Crusader Agreement (GBCA)} (taken from \cite{Abraham2022Efficient,Attiya2023Multi})
Processes start with a binary input $\{0,1\}$ and may return a value in $\{(0,A), (0,B), (\bot,C), (1,B), (1,A)\}$.
A program implements GBCA if it satisfies four properties: Validity, Graded Agreement, Binding, Termination.
\begin{itemize}
    \item Validity: if a correct process does not return $(b,A)$ then a correct process had $1-b$ as input.
    \item Graded Agreement: if two correct processes return $(b,X)$ and $(b',X')$ then $b = 0 \implies b' \ne 1$ and $X = A \implies X' \ne \bot$.
    \item Binding: when a correct process returns then a bound value $b$ is defined and, in all extensions of the execution, correct processes do not return $(1-b,B)$ or $(1-b,A)$.
    \item Termination: if $n-f$ correct processes take part to the protocol then all correct processes will eventually return.
\end{itemize}

\begin{protocol}\label{prot:SpecGBCA}
    \lean{GBCA.Spec}
    %\leanok
    Blablabla
    \begin{algorithm}
        \caption{Ideal Implementation of GBCA}
        \begin{algorithmic}
            \State \textbf{Input:} $b \in \{0,1\}$
            \medskip
            \State \textbf{Code for every process:}
            \State send $\langle input, b\rangle$ to the third party 
            \State \textbf{upon} receiving $\langle output, (b,g)\rangle$ from third party 
            \State \quad return $(b,g)$
            \medskip
            \State \textbf{Code for third party:}
            \State $bind,grade,call \gets \bot,\bot,(\bot)^{Id}$
            \State \textbf{upon} receiving $\langle input, b\rangle$ from a process $id$:
            \State \quad $call(id) \gets b$
            \State \textbf{upon} $|F| + \{id \in Id \setminus F \mid call(id) = b\} \ge n-f$:
            \State \quad $bind \gets b$
            \State \quad \textbf{if} $\{id \in Id \setminus F \mid call(id) = 1-b\} = 0$:
            \State \quad\quad $grade \gets A$
            \State \quad \textbf{else}
            \State \quad\quad $grade \gets C$
            \State \textbf{upon} $bind = \bot$ and $grade = A$ and not sent $\langle output, \cdot \rangle$ to process $id$:
            \State \quad send $\langle bind, B\rangle$ or $\langle output, (bind, A)\rangle$ to process $id$
            \State \textbf{upon} $bind = \bot$ and $grade = C$ and not sent $\langle output, \cdot \rangle$ to process $id$:
            \State \quad send $\langle bind, B\rangle$ or $\langle output, (\bot, C)\rangle$ to process $id$
        \end{algorithmic}
    \end{algorithm}
\end{protocol}

\begin{protocol}\label{prot:ImplGBCA}
    \lean{GBCA.Impl}
    %\leanok
    Parameters: $n$ processes with at most $f < n/3$ corrupted. Each process $p_i$ has input $b_i \in \{0,1\}$.
    \begin{algorithm}
        \caption{Implementation of GBCA taken from \cite{Abraham2022Efficient}}
        \begin{algorithmic}
            \State \textbf{Input:} $b \in \{0,1\}$
            \medskip
            \State send $\langle \textsc{input}, b_i \rangle$ to all processes
            \medskip
            \State \textbf{upon} receiving $\langle \textsc{input}, b \rangle$ from at least $f + 1$ processes and not sent $\langle \textsc{input}, b\rangle$ yet:
            \State \quad send $\langle \textsc{input}, b\rangle$ to all
            \medskip
            \State \textbf{upon} receiving $\langle \textsc{input}, b \rangle$ from at least $n - f$ processes:
            \State \quad $Approved \gets Approved \cup \{b\}$
            \State \quad \textbf{if} not sent $\langle \textsc{echo}, \cdot \rangle$ yet \textbf{then} send $\langle \textsc{echo}, b\rangle$ to all
            \medskip
            \State wait until 
            \State \quad (a) receiving $\langle \textsc{echo}, b \rangle$ from at least $n - f$ processes \textbf{or}
            \State \quad (b) receiving $\langle \textsc{echo}, \cdot \rangle$ from at least $n - f$ processes and $|Approved| > 1$
            \State \quad \textbf{if} (a) \textbf{then} send $\langle vote, b\rangle$ to all \textbf{else} send $\langle vote, \bot\rangle$ to all
            \medskip
            \State wait until 
            \State \quad (a) receiving $\langle \textsc{vote}, b \rangle$ from at least $n - f$ processes \textbf{or}
            \State \quad (b) receiving $\langle \textsc{vote}, \cdot \rangle$ from at least $n - f$ processes and $|Approved| > 1$
            \State \quad \textbf{if} (a) \textbf{then} send $\langle bind, b\rangle$ to all \textbf{else} send $\langle bind, \bot\rangle$ to all
            \medskip
            \State wait until 
            \State \quad (a) receiving $\langle \textsc{vote}, b \rangle$ from at least $n - f$ processes \textbf{or}
            \State \quad (b) receiving $\langle \textsc{vote}, \cdot \rangle$ from at least $n - f$ processes and there exists $b \in \{0,1\}$ such that $\langle bind, b\rangle$ has been received once and $\langle vote, b\rangle$ has been received from $t+1$ processes and $|Approved| > 1$ \textbf{or}
            \State \quad (c) receiving $\langle \textsc{vote}, \bot \rangle$ from at least $n - f$ processes and $|Approved| > 1$
            \State \quad \textbf{if} (a) \textbf{then} return $(b,A)$ \textbf{elif} (b) \textbf{then} return $(b,B)$ \textbf{else} return $(\bot, C)$
        \end{algorithmic}
    \end{algorithm}
\end{protocol}

\paragraph{Weak Common Coin (WCC)}
Processes start with no input and may return a bit.
A program implements WCC if it satisfies three properties: $\epsilon$-Goodness, Unpredictability, Termination.
\begin{itemize}
    \item $\epsilon$-Goodness: with probability $\epsilon$, all correct processes return 0 (resp. 1).
    \item Unpredictability: Until $t+1$ processes have called WCC, all outcomes are still possible.
    \item Termination: if $n-f$ correct processes take part to the protocol then all correct processes will eventually return.
\end{itemize}

\begin{protocol}\label{prot:SpecWCC}
    \lean{WCC.Spec}
    %\leanok
    Blablabla
    \begin{algorithm}
        \caption{Ideal Implementation of WCC}
        \begin{algorithmic}
            \State \textbf{Input:} none
            \medskip
            \State \textbf{Code for every process:}
            \State send $\langle input, b\rangle$ to the third party 
            \State \textbf{upon} receiving $\langle output, b\rangle$ from third party 
            \State \quad return $b$
            \medskip
            \State \textbf{Code for third party:}
            \State $r \gets \mu$ \Comment{$\mu: 1 \mapsto \epsilon; 0 \mapsto \epsilon; \bot \mapsto 1-2\epsilon$}
            \State \textbf{if} $r \ne \bot$:
            \State \quad send $\langle output, r \rangle$ to all
            \State \textbf{alse}
            \State \quad for all $id \in Id$, send $\langle output, 0\rangle$ \textbf{or} $\langle output, 1\rangle$ to process $id$
        \end{algorithmic}
    \end{algorithm}
\end{protocol}

\begin{protocol}\label{prot:ImplWCC}
    \lean{WCC.Impl}
    %\leanok
    Requires Gather and RSD
    \begin{algorithm}
        \caption{Implementation of WCC (yet to be determined)}
        \begin{algorithmic}
            \State $x \gets 0$
            \If{$y > 0$}
                \State loop
            \EndIf
        \end{algorithmic}
    \end{algorithm}
\end{protocol}

\begin{definition}\label{def:WCCConcrete}
    \lean{WCC.Concrete}
    %\leanok
    \uses{prot:ImplWCC, def:GatherConcrete, def:RSDConcrete}
    $\mathsf{WCC.Concrete} := \mathsf{WCC.Impl}[\mathsf{Gather.Concrete}, \mathsf{RSD.Concrete}]$, the instantiation of $\mathsf{WCC.Impl}$ with fully concrete sub-protocols.
\end{definition}

\begin{definition}\label{def:WCCHybrid}
    \lean{WCC.Hybrid}
    %\leanok
    \uses{prot:ImplWCC, prot:SpecGather, prot:SpecRSD}
    $\mathsf{WCC.Hybrid} := \mathsf{WCC.Impl}[\mathsf{Gather.Spec}, \mathsf{RSD.Spec}]$, the instantiation of $\mathsf{WCC.Impl}$ with ideal specifications of the sub-protocols.
\end{definition}

\subsection{WCC Sub-Protocols}

\paragraph{Gather} (taken from \cite{Stern2021Living,SterA2024Gather,Attiya2025Beyond})
Processes start with an input $x \in X$ and may return a partial function $f: Proc \rightharpoonup X$.
A program implements Gather if it satisfies four properties: Validity, Agreement, Common Core, Binding, Termination.
\begin{itemize}
    \item Validity: if a correct process returns $f$ and $f(q)$ is defined and $q$ is correct then $f(q)$ is $q$'s input.
    \item Agreement: if two correct processes return $f$ and $g$ then $f(q)$ and $g(q)$ are equal or one of them is undefined.
    \item Binding Common Core: Once a correct process has returned, there exists $S \subseteq Proc$ of size at least $n-f$ such that all $f$ returned by a correct process in the future is defined on $S$.
    \item Termination: if $n-f$ correct processes take part to Gather then all correct processes will eventually return.
\end{itemize}

\begin{protocol}\label{prot:SpecGather}
    \lean{Gather.Spec}
    %\leanok
    Blablabla
    \begin{algorithm}
        \caption{Ideal Implementation of Gather}
        \begin{algorithmic}
            \State \textbf{Input:} $m$
            \medskip
            \State \textbf{Code for every process:}
            \State send $\langle input, m\rangle$ to the third party 
            \State \textbf{upon} receiving $\langle output, f\rangle$ from third party:
            \State \quad return $f$
            \medskip
            \State \textbf{Code for third party:}
            \State $call, bind \gets \emptyset, \emptyset$
            \State \textbf{upon} receiving $\langle input, m\rangle$ from process $id$: 
            \State \quad $call \gets call \cup \{(id,m)\}$
            \State \textbf{if} $|F \cup \{id \in Id \setminus F \mid call[id] \ne \bot\}| > n-f$ and $bind = \bot$:
            \State \quad $bind \gets S \subseteq F \cup \{id \in Id \setminus F \mid (id,\cdot) \in call\}$ \Comment{Non-deterministic affectation}
            \State \textbf{upon} $bind \ne \bot$:
            \State \quad for all $id \in Id$, 
                            select $bind \subseteq S \subseteq F \cup \{id \in Id \setminus F \mid (id,\cdot) \in call\}$ and \\
                            send $\langle output, \{(id,m) \in call \mid id \in S \setminus F\} \cup \{(id, \cdot) \mid id \in S \cap F\}\rangle$ to process $id$
        \end{algorithmic}
    \end{algorithm}
\end{protocol}

\begin{protocol}\label{prot:ImplGather}
    \lean{Gather.Impl}
    %\leanok
    Might require BRB 
    \begin{algorithm}
        \caption{Implementation of Gather from \cite{Attiya2025Beyond}}
        \begin{algorithmic}
            \State \textbf{Input:} $b \in \{0,1\}$
            \State \textbf{Sub-Protocols:} $|Id|$ instances of BRB (process $id$ is leader in $BRB_{id}$)
            \medskip
            \State $BRB_{id}.call(m)$
            \medskip
            \State \textbf{upon} $BRB_i.return(m')$:
            \State \quad $AP \gets AP \cup \{(i,m')\}$
            \medskip
            \State \textbf{upon} $|AP| > n-f$ and not sent $\langle \textsc{echo}, \cdot\rangle$:
            \State \quad send $\langle echo, AP\rangle$ to all
            \medskip
            \State \textbf{upon} receiving approved $\langle \textsc{echo}, AP_{id} \rangle$ from at least $n - f$ processes and not sent $\langle \textsc{vote}, \cdot\rangle$:
            \State \quad send $\langle \textsc{vote}, \bigcup AP_{id} \rangle$
            \medskip
            \State \textbf{upon} receiving approved $\langle \textsc{vote}, AP_{id} \rangle$ from at least $n - f$ processes and not sent $\langle \textsc{bind}, \cdot\rangle$:
            \State \quad send $\langle \textsc{bind}, \bigcup AP_{id} \rangle$
            \medskip
            \State \textbf{upon} receiving approved $\langle \textsc{bind}, AP_{id} \rangle$ from at least $n - f$ processes:
            \State \quad return $\bigcup AP_{id}$
        \end{algorithmic}
    \end{algorithm}
\end{protocol}

\begin{definition}\label{def:GatherConcrete}
    \lean{Gather.Concrete}
    %\leanok
    \uses{prot:ImplGather, prot:ImplBRB}
    $\mathsf{Gather.Concrete} := \mathsf{Gather.Impl}[\mathsf{BRB.Impl}]$, the instantiation of $\mathsf{Gather.Impl}$ with the concrete implementation of BRB.
\end{definition}

\begin{definition}\label{def:GatherHybrid}
    \lean{Gather.Hybrid}
    %\leanok
    \uses{prot:ImplGather, prot:SpecBRB}
    $\mathsf{Gather.Hybrid} := \mathsf{Gather.Impl}[\mathsf{BRB.Spec}]$, the instantiation of $\mathsf{Gather.Impl}$ with the ideal specification of BRB.
\end{definition}

\paragraph{Random Secret Draw (RSD)}
The RSD protocol is composed of two operations: Assignment and Retrieval. 
Assignment affects a random and secret value to each process and Retrieval provides these values to processes.
Moreover, the API of RSD provides a notification function for each process that records which processes have a value assigned.
\begin{itemize}
    \item Assignment Validity: if all correct processes call Assignment, then all correct processes will eventually be receive a notification for each correct process.
    \item Assignment Totality: if a correct process receives a notification for process $id$ then all correct processes will get such a notification.
    \item Randomness: If no correct process has called Assignment then the value returned by Retrieval for any process follows the uniform distribution.
    \item Secrecy: If no correct process has called Retrieval then the adversary has no information about the return value of Retrieval.
    \item Retrieval Termination: if all correct processes call Retrieval then the operation terminates for all correct processes 
\end{itemize}

\begin{protocol}\label{prot:SpecRSD}
    \lean{RSD.Spec}
    %\leanok
    Blablabla
    \begin{algorithm}
        \caption{Ideal Implementation of RSD}
        \begin{algorithmic}
            \State $x \gets 0$
            \If{$y > 0$}
                \State loop
            \EndIf
        \end{algorithmic}
    \end{algorithm}
\end{protocol}

\begin{protocol}\label{prot:ImplRSD}
    \lean{RSD.Impl}
    %\leanok
    Require AVSS and BRB
    \begin{algorithm}
        \caption{Implementation of RSD from \cite{Freitas2022Distributed}}
        \begin{algorithmic}
            \State \textbf{Input:} $b \in \{0,1\}$
            \State \textbf{Sub-Protocols:} $|Id|$ instances of BRB (process $id$ is leader in $BRB_{id}$) and $|Id|^2$ instances of AVSS (process $id$ is dealer in $AVSS_{id}^{k}$)
            \medskip
            \State \textbf{Operation:} $RSD.Sh()$
            \State $retrieveEnabled, S \gets \bot, \emptyset$
            \State for all $id' \in Id$, $x_{id'} \gets \mu$ \Comment{$\mu$: Uniform distribution over the domain of secerts}
            \State $AVSS^{id'}_{id}.Sh.call(x_{id'})$
            \State \textbf{upon} $AVSS_{id'}^{id}.Sh.return()$:
            \State \quad $S \gets S \cup \{id'\}$
            \State \textbf{upon} $|AP| > n-f$ and not called $BRB_{id}$:
            \State \quad $BRB_{id}.call(S)$
            \State \textbf{upon} $BRB_{id'}.return(S')$ and $\forall k \in S', AVSS^{id'}_k.Sh.return()$:
            \State \quad $src[id'] \gets S'$
            \State \quad \textbf{event} $ValueAssigned(id')$
            \State \quad %\textbf{if} $retrieveEnabled$ \textbf{then} 
            $\forall k \in S', AVSS^{id'}_k.Rec.call()$
            \medskip
            %\State \textbf{Operation:} $RSD.EnableRec()$
            %\State $retrieveEnabled \gets \top$
            %\State $\forall id' \in Id$ \textbf{if} $ValueAssigned(id')$ \textbf{then} $AVSS.Rec.call()$
            %\medskip
            \State \textbf{Operation:} $RSD.Rec()$
            \State return $id' \mapsto \sum_{k \in src[id']}AVSS_k^{id'}.Rec.return()$
        \end{algorithmic}
    \end{algorithm}
\end{protocol}

\begin{definition}\label{def:RSDConcrete}
    \lean{RSD.Concrete}
    %\leanok
    \uses{prot:ImplRSD, prot:ImplBRB, prot:SpecAVSS}
    $\mathsf{RSD.Concrete} := \mathsf{RSD.Impl}[\mathsf{BRB.Impl}, \mathsf{AVSS.Spec}]$, the instantiation of $\mathsf{RSD.Impl}$ with the concrete implementation of BRB. $\mathsf{AVSS}$ is used as a black-box, so its specification is plugged in directly.
\end{definition}

\begin{definition}\label{def:RSDHybrid}
    \lean{RSD.Hybrid}
    %\leanok
    \uses{prot:ImplRSD, prot:SpecBRB, prot:SpecAVSS}
    $\mathsf{RSD.Hybrid} := \mathsf{RSD.Impl}[\mathsf{BRB.Spec}, \mathsf{AVSS.Spec}]$, the instantiation of $\mathsf{RSD.Impl}$ with the ideal specifications of BRB and AVSS.
\end{definition}

\subsection{Basic Primitives}

\paragraph{Byzantine Reliable Broadcast (BRB)} (taken from \cite{Bracha1987Asynchronous})
Each instance of BRB has a designated \emph{leader} with an input message $m$ and processes may return a received message.
A program implements BRB if it satisfies three properties: Validity, Totality and Agreement.
\begin{itemize}
    \item Validity: if the leader is not corrupted then all correct processes must return its input.
    \item Agreement: if two correct processes return then they return the same value.
    \item Totality: if a correct process returns then all correct processes eventually return.
\end{itemize}

\begin{protocol}\label{prot:SpecBRB}
    \lean{BRB.Spec}
    %\leanok
    Blablabla
    \begin{algorithm}
        \caption{Ideal Implementation of BRB}
        \begin{algorithmic}
            \State \textbf{Input:} $m \in \{0,1\}^*$ for the leader
            \medskip
            \State \textbf{Code for the leader:}
            \State send $\langle input, m\rangle$ to the third party 
            \medskip
            \State \textbf{Code for every process:}
            \State \textbf{upon} receiving $\langle output, m\rangle$ from third party:
            \State \quad return $m$
            \medskip
            \State \textbf{Code for third party:}
            \State \textbf{upon} receiving $\langle input, m\rangle$ from leader:
            \State \quad \textbf{if} leader is correct:
            \State \quad\quad send $\langle output, m\rangle$ to all
            \State \quad \textbf{else}
            \State \quad\quad send $\langle output, m'\rangle$ to all \Comment{Arbitrary $m'$}
        \end{algorithmic}
    \end{algorithm}
\end{protocol}

\begin{protocol}\label{prot:ImplBRB}
    \lean{BRB.Impl}
    %\leanok
    Requires no sub-protocol. Parameters: $n$ processes with at most $f < n/3$ corrupted; leader $\ell$ with input $m$.
    \begin{algorithm}
        \caption{Implementation of BRB taken from \cite{Bracha1987Asynchronous}}
        \begin{algorithmic}
            \State \textbf{Code for the leader $\ell$ with input $m$:}
            \State \quad send $\langle \textsc{init}, m \rangle$ to all processes
            \medskip
            \State \textbf{Code for every process $p_i$:}
            \State \textbf{upon} first reception of $\langle \textsc{init}, m \rangle$ from $\ell$ \textbf{do}
                \State \quad send $\langle \textsc{echo}, m \rangle$ to all processes
            \State \textbf{upon} reception of $\langle \textsc{echo}, m \rangle$ from at least $n-f$ distinct processes \textbf{do}
                \State \quad \textbf{if} $\langle \textsc{ready}, \cdot \rangle$ has not yet been sent \textbf{then}
                    \State \quad\quad send $\langle \textsc{ready}, m \rangle$ to all processes
            \State \textbf{upon} reception of $\langle \textsc{ready}, m \rangle$ from at least $f+1$ distinct processes \textbf{do}
                \State \quad \textbf{if} $\langle \textsc{ready}, \cdot \rangle$ has not yet been sent \textbf{then}
                    \State \quad\quad send $\langle \textsc{ready}, m \rangle$ to all processes
            \State \textbf{upon} reception of $\langle \textsc{ready}, m \rangle$ from at least $2f+1$ distinct processes \textbf{do}
                \State \quad deliver $m$
        \end{algorithmic}
    \end{algorithm}
\end{protocol}


\paragraph{Asynchronous Verifiable Secret Sharing (AVSS)}
AVSS is actually a pair of protocol $(Sh, Ret)$ with a designated \emph{dealer} owning a secret $s$.
A program implements a $(1-\epsilon)$-correct AVSS if it satisfies six properties: Valid Termination, Total Termination, $Rec$ Termination, Binding, Secrecy.
\begin{itemize}
    \item Validity: if the dealer is correct and a correct process returns from $Ret$ then the return value is the dealer's secret.
    \item Sharing Termination: if the dealer is correct and calls $Sh$ then all correct processes will enventually return from $Sh$.
    \item Sharing Totality: if a correct process returns from $Sh$ then all correct processes will eventually return from $Sh$.
    \item Retrieval Termination: if all correct processes call $Ret$ then all correct processes will eventually return from $Ret$.
    \item Commitment: if a correct process returns from $Sh$ then there exists a secret $s$ such that correct process can only retrieve $s$ from that point.
    \item Secrecy: if the dealer is correct and no correct process has called $Ret$ then the adversary has no information about the future retrived value.
\end{itemize}

\begin{protocol}\label{prot:SpecAVSS}
    \lean{AVSS.Spec}
    %\leanok
    Blablabla
    \begin{algorithm}
        \caption{AVSS Specification}
        \begin{algorithmic}
            \State $x \gets 0$
            \If{$y > 0$}
                \State loop
            \EndIf
        \end{algorithmic}
    \end{algorithm}
\end{protocol}

AVSS is used as a blackbox so no implementation.

\subsection{Summary of Properties}

The linear safety properties are preserved by Segala's probabilistic forward simulation. 
The branching safety properties can be transformed into linear ones by enhancing the 
interface with an external binding label.
Segala's simulation does not preserve liveness by default. 
However, it can be strengthen following \cite{Dongol2022Weak} to preserve liveness 
in a complete manner.

\begin{table}[ht]
    \centering
    \begin{tabular}{|c||c|c|c|c|c|c|}
        \hline
        & Linear Safety & Branching Safety & Liveness \\
        \hline
        ABA & Validity \& Agreement & &  \\
        GBCA & Validity \& Agreement & Binding & Termination \\
        WCC & & & Termination  \\
        Gather & Validity \& Agreement & Binding Common Core & Termination  \\
        RSD & & & \makecell{Assignment Validity \& \\ Assignment Totality \& \\ Retrieval Termination} \\
        BRB & Validity \& Agreement & & Totality \\
        AVSS & Validity & Commitment & \makecell{Sharing Termination \& \\ Sharing Totality \& \\ Retrieval Termination} \\
        \hline
    \end{tabular}
    \caption{Table of CTL Properties}
\end{table}

Segala's simulation was designed to be sound and complete for trace distribution so
it preserves probabilistic safety propreties.
Note that Unpredictability is a probabilistic safety property.
For probabilistic liveness, we will rely on the liveness preservation as it is sound 
(but not complete). 
The lack of completeness is not an isue because the limitations of the termination
proof rule is probabilistic convergence which is never dealt with by our simulations.
The probabilistic convergence argument will only be necessary to prove that the ideal 
implementation of ABA almost-surely terminates and can be discharged to the 
super-Martingale proof rule of \cite{Enea2026Verifying}.
For Secrecy, we use the same trick as for branching: enhancing the interface. 


\begin{table}[ht]
    \centering
    \begin{tabular}{|c||c|c|c|c|c|c|}
        \hline
        & Probabilistic Safety & Probabilistic Liveness & Secrecy \\
        \hline
        ABA & & Almost-Sure Termination & \\
        GBCA & & & \\
        WCC & $\epsilon$-Goodness \& Unpredictability & & \\
        Gather & & & \\
        RSD & Randomness & & Secrecy \\
        BRB & & & \\
        AVSS & & & Secrecy \\
        \hline
    \end{tabular}
    \caption{Table of non-CTL Properties}
\end{table}




\section{Proof Decomposition}

The end goal is the correctness of $\mathsf{ABA.Impl}$.
We establish it via a simulation by $\mathsf{ABA.Spec}$ and transport the spec's correctness.
The simulation itself is decomposed recursively along the sub-protocol structure.

For a protocol $P$ with concrete sub-protocols $Q_1,\dots,Q_k$, the step ``$\mathsf{P.Spec}$ simulates $\mathsf{P.Impl}$'' is split into:
\begin{enumerate}
    \item (Substitution) $\mathsf{P.Impl}[\mathsf{Q_1.Spec},\dots,\mathsf{Q_k.Spec}]$ is simulated by $\mathsf{P.Impl}$, using the child simulations $\mathsf{Q_i.Spec}$ simulates $\mathsf{Q_i.Impl}$ and compositionality.
    \item (Core) $\mathsf{P.Spec}$ simulates the hybrid $\mathsf{P.Impl}[\mathsf{Q_1.Spec},\dots,\mathsf{Q_k.Spec}]$.
    \item (Transitivity) Compose (1) and (2).
\end{enumerate}
Leaves of the recursion (no sub-protocol, or only the black-box $\mathsf{AVSS.Spec}$) skip the substitution step.

\subsection{Top-level}

\begin{theorem}\label{thm:ABACorrect}
    \lean{ABA.Correct}
    %\leanok
    \uses{def:ABAConcrete, prot:SpecABA}
    $\mathsf{ABA.Concrete}$ satisfies Validity, Agreement and Almost-Sure Termination.
\end{theorem}
\begin{proof}
    \uses{thm:ABASim, thm:pfs-sound}
    Soundness of probabilistic forward simulation (Theorem~\ref{thm:pfs-sound}) applied to \ref{thm:ABASim} yields $\mathsf{ABA.Concrete} \sqsubseteq_{\mathrm{TD}} \mathsf{ABA.Spec}$, which transports the safety properties (Validity and Agreement) of $\mathsf{ABA.Spec}$ to $\mathsf{ABA.Concrete}$. Almost-Sure Termination is a probabilistic liveness property and is handled separately via the liveness-preserving strengthening of $\preceq_{\mathrm{PFS}}$.
\end{proof}

\subsection{ABA}

\begin{theorem}\label{thm:ABASim}
    \lean{ABA.Sim}
    %\leanok
    \uses{def:ABAConcrete, prot:SpecABA}
    $\mathsf{ABA.Spec}$ simulates $\mathsf{ABA.Concrete}$.
\end{theorem}
\begin{proof}
    \uses{lem:ABASubst, lem:ABACore, thm:pfs-largest-precongruence}
    Probabilistic forward simulation is in particular a pre-order (Theorem~\ref{thm:pfs-largest-precongruence}); chain \ref{lem:ABACore} with \ref{lem:ABASubst} by transitivity.
\end{proof}

\begin{lemma}\label{lem:ABASubst}
    \lean{ABA.SubstSim}
    %\leanok
    \uses{def:ABAConcrete, def:ABAHybrid}
    $\mathsf{ABA.Hybrid}$ is simulated by $\mathsf{ABA.Concrete}$.
\end{lemma}
\begin{proof}
    \uses{thm:GBCASim, thm:WCCSim, thm:pfs-precongruence}
    $\mathsf{ABA.Concrete}$ and $\mathsf{ABA.Hybrid}$ share the same skeleton $\mathsf{ABA.Impl}$ and differ only in their sub-protocol instances. Apply the pre-congruence of $\preceq_{\mathrm{PFS}}$ for parallel composition (Theorem~\ref{thm:pfs-precongruence}) to the child simulations \ref{thm:GBCASim} and \ref{thm:WCCSim}, lifted along the $\mathsf{ABA.Impl}$ context.
\end{proof}

\begin{lemma}\label{lem:ABACore}
    \lean{ABA.CoreSim}
    %\leanok
    \uses{def:ABAHybrid, prot:SpecABA}
    $\mathsf{ABA.Spec}$ simulates $\mathsf{ABA.Hybrid}$.
\end{lemma}

\subsection{WCC}

\begin{theorem}\label{thm:WCCSim}
    \lean{WCC.Sim}
    %\leanok
    \uses{def:WCCConcrete, prot:SpecWCC}
    $\mathsf{WCC.Spec}$ simulates $\mathsf{WCC.Concrete}$.
\end{theorem}
\begin{proof}
    \uses{lem:WCCSubst, lem:WCCCore, thm:pfs-largest-precongruence}
    Probabilistic forward simulation is in particular a pre-order (Theorem~\ref{thm:pfs-largest-precongruence}); chain \ref{lem:WCCCore} with \ref{lem:WCCSubst} by transitivity.
\end{proof}

\begin{lemma}\label{lem:WCCSubst}
    \lean{WCC.SubstSim}
    %\leanok
    \uses{def:WCCConcrete, def:WCCHybrid}
    $\mathsf{WCC.Hybrid}$ is simulated by $\mathsf{WCC.Concrete}$.
\end{lemma}
\begin{proof}
    \uses{thm:GatherSim, thm:RSDSim, thm:pfs-precongruence}
    $\mathsf{WCC.Concrete}$ and $\mathsf{WCC.Hybrid}$ share the same skeleton $\mathsf{WCC.Impl}$ and differ only in the instances of Gather and RSD. Apply the pre-congruence of $\preceq_{\mathrm{PFS}}$ for parallel composition (Theorem~\ref{thm:pfs-precongruence}) to the child simulations \ref{thm:GatherSim} and \ref{thm:RSDSim}, lifted along the $\mathsf{WCC.Impl}$ context.
\end{proof}

\begin{lemma}\label{lem:WCCCore}
    \lean{WCC.CoreSim}
    %\leanok
    \uses{def:WCCHybrid, prot:SpecWCC}
    $\mathsf{WCC.Spec}$ simulates $\mathsf{WCC.Hybrid}$.
\end{lemma}

\subsection{GBCA}

\begin{theorem}\label{thm:GBCASim}
    \lean{GBCA.Sim}
    %\leanok
    \uses{prot:ImplGBCA, prot:SpecGBCA}
    $\mathsf{GBCA.Spec}$ simulates $\mathsf{GBCA.Impl}$. No sub-protocol: direct proof.
\end{theorem}

\subsection{Gather}

\begin{theorem}\label{thm:GatherSim}
    \lean{Gather.Sim}
    %\leanok
    \uses{def:GatherConcrete, prot:SpecGather}
    $\mathsf{Gather.Spec}$ simulates $\mathsf{Gather.Concrete}$.
\end{theorem}
\begin{proof}
    \uses{lem:GatherSubst, lem:GatherCore, thm:pfs-largest-precongruence}
    Probabilistic forward simulation is in particular a pre-order (Theorem~\ref{thm:pfs-largest-precongruence}); chain \ref{lem:GatherCore} with \ref{lem:GatherSubst} by transitivity.
\end{proof}

\begin{lemma}\label{lem:GatherSubst}
    \lean{Gather.SubstSim}
    %\leanok
    \uses{def:GatherConcrete, def:GatherHybrid}
    $\mathsf{Gather.Hybrid}$ is simulated by $\mathsf{Gather.Concrete}$.
\end{lemma}
\begin{proof}
    \uses{thm:BRBSim, thm:pfs-precongruence}
    $\mathsf{Gather.Concrete}$ and $\mathsf{Gather.Hybrid}$ share the same skeleton $\mathsf{Gather.Impl}$ and differ only in the BRB instances. Apply the pre-congruence of $\preceq_{\mathrm{PFS}}$ for parallel composition (Theorem~\ref{thm:pfs-precongruence}) to the child simulation \ref{thm:BRBSim}, lifted along the $\mathsf{Gather.Impl}$ context (one instance per leader).
\end{proof}

\begin{lemma}\label{lem:GatherCore}
    \lean{Gather.CoreSim}
    %\leanok
    \uses{def:GatherHybrid, prot:SpecGather}
    $\mathsf{Gather.Spec}$ simulates $\mathsf{Gather.Hybrid}$.
\end{lemma}

\subsection{RSD}

\begin{theorem}\label{thm:RSDSim}
    \lean{RSD.Sim}
    %\leanok
    \uses{def:RSDConcrete, prot:SpecRSD}
    $\mathsf{RSD.Spec}$ simulates $\mathsf{RSD.Concrete}$.
\end{theorem}
\begin{proof}
    \uses{lem:RSDSubst, lem:RSDCore, thm:pfs-largest-precongruence}
    Probabilistic forward simulation is in particular a pre-order (Theorem~\ref{thm:pfs-largest-precongruence}); chain \ref{lem:RSDCore} with \ref{lem:RSDSubst} by transitivity.
\end{proof}

\begin{lemma}\label{lem:RSDSubst}
    \lean{RSD.SubstSim}
    %\leanok
    \uses{def:RSDConcrete, def:RSDHybrid}
    $\mathsf{RSD.Hybrid}$ is simulated by $\mathsf{RSD.Concrete}$. Only BRB is substituted; $\mathsf{AVSS.Spec}$ stays fixed on both sides.
\end{lemma}
\begin{proof}
    \uses{thm:BRBSim, thm:pfs-precongruence}
    $\mathsf{RSD.Concrete}$ and $\mathsf{RSD.Hybrid}$ share the same skeleton $\mathsf{RSD.Impl}$ and the same $\mathsf{AVSS.Spec}$ instances; only the BRB instances are substituted. Apply the pre-congruence of $\preceq_{\mathrm{PFS}}$ for parallel composition (Theorem~\ref{thm:pfs-precongruence}) to the child simulation \ref{thm:BRBSim}, lifted along the $\mathsf{RSD.Impl}$ context (one BRB instance per leader).
\end{proof}

\begin{lemma}\label{lem:RSDCore}
    \lean{RSD.CoreSim}
    %\leanok
    \uses{def:RSDHybrid, prot:SpecRSD}
    $\mathsf{RSD.Spec}$ simulates $\mathsf{RSD.Hybrid}$.
\end{lemma}

\subsection{BRB}

\begin{theorem}\label{thm:BRBSim}
    \lean{BRB.Sim}
    %\leanok
    \uses{prot:ImplBRB, prot:SpecBRB}
    $\mathsf{BRB.Spec}$ simulates $\mathsf{BRB.Impl}$. No sub-protocol: direct proof.
\end{theorem}


\bibliographystyle{alpha}    % or plain, abbrv, alphaurl, etc.
\bibliography{biblio}

\appendix

\clearpage
\section{Implementation's Pseudo-code}

\begin{algorithm}
    \label{alg:ABA}
    \caption{Implementation of \textsf{ABA} from \cite{Abraham2022Efficient}}
    \begin{algorithmic}
        \State \textbf{Input:} $\texttt{b} \in \{0,1\}$
        \State \textbf{Sub-Protocols:} Countably many instances of \textsf{GBCA} and \textsf{WCC}
        \medskip
        \State $\texttt{r} \gets 1$
        \State \textbf{loop}
        \State \quad $(\texttt{b}, \texttt{g}) \gets \textsf{GBCA}_{\texttt{r}}(\texttt{b})$
        \Comment{$\texttt{b} \in \{0,1,\bot\}$, $\texttt{g} \in \{A, B, C\}$}
        \State \quad  $\texttt{c} \gets \textsf{WCC}_{\texttt{r}}()$
        \State \quad \textbf{if} $\texttt{b} = \bot$:
        \State \quad\quad $\texttt{b} \gets \texttt{c}$
        \State \quad \textbf{elif} $\texttt{g} = A$:
        \State \quad\quad send $\langle\textsc{decided}, \texttt{b}\rangle$ to all
        \State \quad $\texttt{r} \gets \texttt{r} + 1$
        \medskip
        \State \textbf{upon} receiving $\langle\textsc{decided}, \texttt{b}\rangle$ from $f+1$
        processes and having not sent it:
        \State \quad send $\langle\textsc{decided}, \texttt{b}\rangle$ to all
        \State \textbf{upon} receiving $\langle\textsc{decided}, \texttt{b}\rangle$ from $n-f$
        processes and having sent it:
        \State \quad return $\texttt{b}$
    \end{algorithmic}
\end{algorithm}

\begin{algorithm}
    \label{alg:GBCA}
    \caption{Implementation of \textsf{GBCA} from \cite{Abraham2022Efficient}}
    \begin{algorithmic}
        \State \textbf{Input:} $\texttt{b} \in \{0,1\}$
        \medskip
        \State send $\langle \textsc{input}, \texttt{b} \rangle$ to all processes
        \State \textbf{upon} receiving $\langle \textsc{input}, \texttt{b} \rangle$ from at least
        $f + 1$ processes and not sent $\langle \textsc{input}, \texttt{b}\rangle$ yet:
        \State \quad send $\langle \textsc{input}, \texttt{b}\rangle$ to all
        \State \textbf{upon} receiving $\langle \textsc{input}, \texttt{b} \rangle$ from at least
        $n - f$ processes:
        \State \quad $\texttt{Valid} \gets \texttt{Valid} \cup \{\texttt{b}\}$
        \State \quad \textbf{if} not sent $\langle \textsc{echo}, \cdot \rangle$ yet \textbf{then}
        send $\langle \textsc{echo}, \texttt{b}\rangle$ to all
        \State \textbf{wait until}
        \State \quad (a) receiving $\langle \textsc{echo}, \texttt{b} \rangle$ from at least
        $n - f$ processes \textbf{or}
        \State \quad (b) receiving $\langle \textsc{echo}, \cdot \rangle$ from at least $n - f$
        processes and $|\texttt{Valid}| > 1$
        \State \quad \textbf{if} (a) \textbf{then} send $\langle \textsc{vote}, \texttt{b}\rangle$
        to all \textbf{else} send $\langle \textsc{vote}, \bot\rangle$ to all
        \State \textbf{wait until}
        \State \quad (a) receiving $\langle \textsc{vote}, \texttt{b} \rangle$ from at least
        $n - f$ processes \textbf{or}
        \State \quad (b) receiving $\langle \textsc{vote}, \cdot \rangle$ from at least $n - f$
        processes and $|\texttt{Valid}| > 1$
        \State \quad \textbf{if} (a) \textbf{then} send $\langle \textsc{bind}, \texttt{b}\rangle$
        to all \textbf{else} send $\langle \textsc{bind}, \bot\rangle$ to all
        \State \textbf{wait until}
        \State \quad (a) receiving $\langle \textsc{bind}, \texttt{b} \rangle$ from at least
        $n - f$ processes \textbf{or}
        \State \quad (b) receiving $\langle \textsc{bind}, \cdot \rangle$ from at least $n - f$
        processes and \\
        \quad  there exists $\texttt{b} \in \{0,1\}$ such that
        $\langle \textsc{bind}, \texttt{b}\rangle$ has been received once and \\
        \quad $\langle \textsc{vote}, \texttt{b}\rangle$ has been received from $f+1$ processes
        and $|\texttt{Valid}| > 1$ \textbf{or}
        \State \quad (c) receiving $\langle \textsc{bind}, \bot \rangle$ from at least $n - f$
        processes and $|\texttt{Valid}| > 1$
        \State \quad \textbf{if} (a) \textbf{then} return $(\texttt{b},A)$ \textbf{elif} (b)
        \textbf{then} return $(\texttt{b},B)$ \textbf{else} return $(\bot, C)$
    \end{algorithmic}
\end{algorithm}

\begin{algorithm}
    \label{alg:WCC}
    \caption{Implementation of \textsf{WCC} for process $p_{id}$ adapted from
    \cite{Freitas2022Distributed}}
    \begin{algorithmic}
        \State \textbf{Input:}
        \State \textbf{Sub-Protocols:} 1 instance of \textsf{Gather}, $n$ instances of
        \textsf{SRSD} (process $p_k$ is the owner of $\textsf{SRSD}_k$)
        \medskip
        \State $\forall k \in [n], \textsf{SRSD}_{k}.\textsf{Asg}.call()$
        \State \textbf{wait until} $\textsf{SRSD}_{id}.\textsf{Asg}.return()$
        \State \quad $\textsf{Gather}.call(id)$
        \State \textbf{upon} $\textsf{Gather}.return(S)$
        \State \quad $\forall k \in S$, $\textsf{SRSD}_k.\textsf{Rec}.call()$
        \State \quad $\texttt{T} \gets \emptyset$
        \State \textbf{upon} $\textsf{SRSD}_k.\textsf{Rec}.return(x)$ and $k \in S$
        \State \quad $\texttt{T} \gets \texttt{T} \cup \{x\}$
        \State \textbf{upon} $|\texttt{T}| = |S|$
        \State \quad return 0 if $0 \in \texttt{T}$ else 1
    \end{algorithmic}
\end{algorithm}

\begin{algorithm}
    \label{alg:Gather}
    \caption{Implementation of \textsf{Gather} from \cite{Attiya2025Beyond}}
    \begin{algorithmic}
        \State \textbf{Input:} $\texttt{m}$
        \State \textbf{Sub-Protocols:} $n$ instances of \textsf{BRB}
        ($p_k$ is leader in $\textsf{BRB}_k$)
        \medskip
        \State $\textsf{BRB}_{id}.call(\texttt{m})$
        \State \textbf{upon} $\textsf{BRB}_k.return(\texttt{m'})$:
        \State \quad $\texttt{AP} \gets \texttt{AP} \cup \{(k, \texttt{m'})\}$
        \State \textbf{upon} $|\texttt{AP}| \ge n-f$ and not sent
        $\langle \textsc{echo}, \cdot\rangle$:
        \State \quad send $\langle \textsc{echo}, \texttt{AP}\rangle$ to all
        \State \textbf{upon} receiving approved $\langle \textsc{echo}, \texttt{AP}_{id} \rangle$
        from at least $n - f$ processes and not sent $\langle \textsc{vote}, \cdot\rangle$:
        \State \quad send $\langle \textsc{vote}, \bigcup_{id} \texttt{AP}_{id} \rangle$ to all
        \State \textbf{upon} receiving approved $\langle \textsc{vote}, \texttt{AP}_{id} \rangle$
        from at least $n - f$ processes and not sent $\langle \textsc{bind}, \cdot\rangle$:
        \State \quad send $\langle \textsc{bind}, \bigcup_{id} \texttt{AP}_{id} \rangle$ to all
        \State \textbf{upon} receiving approved $\langle \textsc{bind}, \texttt{AP}_{id} \rangle$
        from at least $n - f$ processes:
        \State \quad return $\bigcup_{id} \texttt{AP}_{id}$
    \end{algorithmic}
\end{algorithm}

\begin{algorithm}
    \label{alg:SRSD}
    \caption{Implementation of \textsf{SRSD} for process $p_{id}$ adapted from
    \cite{Freitas2022Distributed}}
    \begin{algorithmic}
        \State \textbf{Input:}
        \State \textbf{Sub-Protocols:} 1 instance of \textsf{BRB} where $p_o$ is leader and $n$
        instances of \textsf{AVSS} (process $p_k$ is dealer in $\textsf{AVSS}_k$)
        \medskip
        \State \textbf{Protocol:} $\textsf{SRSD}.\textsf{Asg}()$
        \State $\texttt{S}, \texttt{x} \gets \emptyset, \mu$
        \Comment{$\mu$: uniform distribution over the domain of secrets}
        \State $\textsf{AVSS}_{id}.\textsf{Sh}.call(\texttt{x})$
        \State \textbf{upon} $\textsf{AVSS}_k.\textsf{Sh}.return()$ and $p_{id}$ is the owner:
        \State \quad $\texttt{S} \gets \texttt{S} \cup \{k\}$
        \State \textbf{upon} $|\texttt{S}| \ge n-f$ and $p_{id}$ is the owner:
        \State \quad $\textsf{BRB}.call(S)$
        \State \textbf{upon} $\textsf{BRB}.return(S')$ and
        $\forall k \in S', \textsf{AVSS}_k.\textsf{Sh}.return()$:
        \State \quad $\textsf{src} \gets S'$
        \State \quad return
        \medskip
        \State \textbf{Protocol:} $\textsf{SRSD}.\textsf{Rec}()$
        \State $\texttt{sum} \gets 0$
        \State $\forall k \in \texttt{src}, \textsf{AVSS}_k.\textsf{Rec}.call()$
        \State \textbf{upon} $\textsf{AVSS}_k.\textsf{Rec}.return(x)$ and $k \in \texttt{src}$:
        \State \quad $\texttt{sum} \gets \texttt{sum} + x$
        \State \textbf{wait until} all $\textsf{AVSS}_k.\textsf{Rec}$ have returned for
        $k \in \texttt{src}$:
        \State \quad return $\texttt{sum}$
    \end{algorithmic}
\end{algorithm}

\begin{algorithm}
    \label{alg:BRB}
    \caption{Implementation of \textsf{BRB} for process $p_{id}$ from \cite{Bracha1987Asynchronous}}
    \begin{algorithmic}
        \State \textbf{Input:} $\texttt{m}$
        \medskip
        \State \textbf{if} $p_{id}$ is leader:
        \State \quad send $\langle \textsc{init}, \texttt{m} \rangle$ to all
        \State \textbf{upon} receiving $\langle \textsc{init}, \texttt{m} \rangle$ from leader and
        having not sent $\textsc{echo}$ message:
        \State \quad send $\langle \textsc{echo}, \texttt{m} \rangle$ to all
        \State \textbf{upon} receiving $\langle \textsc{echo}, \texttt{m} \rangle$ from $n-f$
        processes and having not sent $\textsc{vote}$ message:
        \State \quad send $\langle \textsc{vote}, \texttt{m} \rangle$ to all
        \State \textbf{upon} receiving $\langle \textsc{vote}, \texttt{m} \rangle$ from $f+1$
        processes and having not sent $\textsc{vote}$ message:
        \State \quad send $\langle \textsc{vote}, \texttt{m} \rangle$ to all
        \State \textbf{upon} receiving $\langle \textsc{vote}, \texttt{m} \rangle$ from $n-f$
        processes:
        \State \quad return $\texttt{m}$
    \end{algorithmic}
\end{algorithm}

\clearpage
\section{Specification}

\begin{transys}
    \label{PLTS:ABA}
    \caption{Encoding of \textsf{ABA}}
    \begin{algorithmic}
        \State \textbf{State:} $call \in \{0,1,\bot\}^{[n]}; ret \in \bool^{[n]};
        bind, val \in \{0,1,\bot\}; coin \in \{0,1,\bot,\top\}; F \subseteq [n]$
        \State \textbf{Label:} $\{\text{call}(id,b), \text{return}(id,b), \text{fail}(id) \mid
        id \in [n], b \in \bool\}$
        \State \textbf{Input:} $\{\text{call}(id,b), \text{fail}(id) \mid id \in [n], b \in \bool\}$
        \State \textbf{Fairness:} all transitions are fair except those labelled by fail and the
        call loops
        \State \textbf{Initial:} $call = \bot^{[n]}; ret = \bot^{[n]}; bind = \bot; out = \bot;
        val = \bot; F = \emptyset$
        \medskip
        \State
        $call[id] = \bot \wedge bind = \bot
        \xrightarrow{\text{call}(id,b)}
        call[id] = b$
        \State
        $\top
        \xrightarrow{\text{call}(id,b)}
        \top$ \Comment{loop for input-enabledness}
        \State
        $|\{id \in [n] \setminus F \mid call[id] \ne \bot\} \cup F| \ge n-f
        \wedge \forall id \in [n] \setminus F, call[id] \ne b
        \xrightarrow{\tau}
        call = \bot^{[n]} \wedge bind = 1-b \wedge val = 1-b \wedge coin = \bot$
        \State
        $|\{id \in [n] \setminus F \mid call[id] \ne \bot\} \cup F| \ge n-f
        \wedge \exists id, id' \in [n] \setminus F, call[id] = 1 \wedge call[id'] = 0
        \xrightarrow{\tau}
        call = \bot^{[n]} \wedge bind = 1-b \wedge coin = \bot$
        \State
        $call = \bot^{[n]} \wedge bind \ne \bot
        \xrightarrow{\tau}
        coin = 1 \text{ with probability } \epsilon; 0 \text{ with probability } \epsilon;
        \top \text{ otherwise}$
        \State
        $call[id] = \bot \wedge bind = coin
        \xrightarrow{\tau}
        call[id] = bind$
        \State
        $call[id] = \bot \wedge bind \ne coin
        \xrightarrow{\tau}
        call[id] = b$ \Comment{$b \in \bool$}
        \State
        $val = b \wedge ret[id] = \bot
        \xrightarrow{\text{return}(id,b)}
        ret[id] = \top$
        \State
        $id \not\in F \wedge |F| < f
        \xrightarrow{\text{fail}(id)}
        F = F \cup \{id\}$
        \State
        $\top
        \xrightarrow{\text{fail}(id)}
        \top$ \Comment{loop for input-enabledness}
    \end{algorithmic}
\end{transys}

\begin{transys}
    \label{PLTS:GBCA}
    \caption{Encoding of \textsf{GBCA}}
    \begin{algorithmic}
        \State \textbf{State:} $call \in \{0,1,\bot\}^{[n]}; ret \in \bool^{[n]};
        bind, grade \in \{0,1,\bot\}; F \subseteq [n]$
        \State \textbf{Label:} $\{\text{call}(id,b), \text{return}(id,x,b), \text{fail}(id) \mid
        id \in [n], x \in (\bool \times \{A,B\})\cup\{(\bot,C)\}, b \in \bool\}$
        \State \textbf{Input:} $\{\text{call}(id,b), \text{fail}(id) \mid id \in [n], b \in \bool\}$
        \State \textbf{Fairness:} all transitions are fair except those labelled by fail and the
        call loops
        \State \textbf{Initial:} $call = \bot^{[n]}; ret = \bot^{[n]}; bind = \bot; grade = \bot;
        F = \emptyset$
        \medskip
        \State
        $call[id] = \bot
        \xrightarrow{\text{call}(id,b)}
        call[id] = b$
        \State
        $\top
        \xrightarrow{\text{call}(id,b)}
        \top$ \Comment{loop for input-enabledness}
        \State
        $|\{id \in [n] \setminus F \mid call[id] \ne \bot\} \cup F| \ge n-f
        \wedge \exists id \in [n] \setminus F, call[id] = b
        \wedge bind = \bot
        \xrightarrow{\tau}
        bind = b$
        \State
        $bind \ne \bot
        \wedge \exists id \in [n] \setminus F, call[id] = 1-bind
        \wedge ret[id] = \bot
        \xrightarrow{\text{return}(id,(bind,B),bind)}
        ret[id] = \top$
        \State
        $bind \ne \bot \wedge grade \in \{\bot,1\} \wedge ret[id] = \bot
        \xrightarrow{\text{return(id,(bind,A),bind)}}
        grade = 1 \wedge ret[id] = \top$
        \State
        $bind \ne \bot
        \wedge \exists id \in [n] \setminus F, call[id] = 1-bind
        \wedge grade \in \{\bot,0\} \wedge ret[id] = \bot
        \xrightarrow{\text{return}(id,(\bot,C),bind)}
        grade = 0 \wedge ret[id] = \top$
        \State
        $id \not\in F \wedge |F| < f
        \xrightarrow{\text{fail}(id)}
        F = F \cup \{id\}$
        \State
        $\top
        \xrightarrow{\text{fail}(id)}
        \top$ \Comment{loop for input-enabledness}
    \end{algorithmic}
\end{transys}

\begin{transys}
    \label{PLTS:WCC}
    \caption{Encoding of \textsf{WCC}}
    \begin{algorithmic}
        \State \textbf{State:} $call, ret \in \bool^{[n]}; val \in \{0,1,\bot,\top\};
        guess \in \{0,1,\bot\}, F \subseteq [n]$
        \State \textbf{Label:} $\{\text{call}(id), \text{return}(id,b), \text{fail}(id),
        \text{guess}(b) \mid id \in [n], b \in \bool\}$
        \State \textbf{Input:} $\{\text{call}(id), \text{fail}(id) \mid id \in [n]\}$
        \State \textbf{Fairness:} all transitions are fair except those labelled by fail and the
        call loops
        \State \textbf{Initial:} $call = \bot^{[n]}; ret = \bot^{[n]}; val = \bot; F = \emptyset$
        \medskip
        \State
        $call[id] = \bot
        \xrightarrow{\text{call}(id)}
        call[id] = \top$
        \State
        $\top
        \xrightarrow{\text{call}(id)}
        \top$ \Comment{loop for input-enabledness}
        \State
        $|\{id \in [n] \setminus F \mid call[id] \ne \bot\} \cup F| > f
        \wedge val = \bot
        \xrightarrow{\tau}
        val = 0$ with probability $\epsilon$; 1 with probability $\epsilon$; $\top$ otherwise
        \State
        $val \in \{b,\top\}
        \wedge ret[id] = \bot
        \xrightarrow{\text{return}(id,b)}
        ret[id] = \top$ \Comment{$b \in \bool$}
        \State
        $id \not\in F \wedge |F| < f
        \xrightarrow{\text{fail}(id)}
        F = F \cup \{id\}$
        \State
        $\top
        \xrightarrow{\text{fail}(id)}
        \top$ \Comment{loop for input-enabledness}
        \State
        $val = \bot \wedge guess = \bot
        \xrightarrow{\text{guess}(b)}
        guess = b$
    \end{algorithmic}
\end{transys}

\begin{transys}
    \label{PLTS:Gather}
    \caption{Encoding of \textsf{Gather}}
    \begin{algorithmic}
        \State \textbf{State:} $call \in (X\cup\{\bot\})^{[n]}; ret \in \bool^{[n]};
        bind, F \subseteq [n]$
        \State \textbf{Label:} $\{\text{call}(id,x), \text{return}(id,g,bind), \text{fail}(id) \mid
        id \in [n], x \in X, g \in [n] \rightharpoonup X, bind \subseteq [n]\}$
        \State \textbf{Input:} $\{\text{call}(id,x), \text{fail}(id) \mid id \in [n], x \in X\}$
        \State \textbf{Fairness:} all transitions are fair except those labelled by fail and the
        call loops
        \State \textbf{Initial:} $call = \bot^{[n]}; ret = \bot^{[n]}; bind = \emptyset;
        F = \emptyset$
        \medskip
        \State
        $call[id] = \bot
        \xrightarrow{\text{call}(id,x)}
        call[id] = x$
        \State
        $\top
        \xrightarrow{\text{call}(id,x)}
        \top$ \Comment{loop for input-enabledness}
        \State
        $call[id] = \bot \wedge id \in F
        \xrightarrow{\tau}
        call[id] = x$
        \State
        $|\{id \in [n] \mid call[id] \ne \bot\}| > n-f
        \wedge bind = \emptyset
        \wedge S \subseteq \{id \in [n] \mid call[id] \ne \bot\}
        \xrightarrow{\tau}
        bind = S$
        \State
        $bind \ne \bot
        \wedge ret[id] = \bot
        \wedge g \text{ is defined on } bind
        \wedge g \text{ agrees with } call
        \xrightarrow{\text{return}(id,g,bind)}
        ret[id] = \top$
        \State
        $id \not\in F \wedge |F| < f
        \xrightarrow{\text{fail}(id)}
        F = F \cup \{id\}$
        \State
        $\top
        \xrightarrow{\text{fail}(id)}
        \top$ \Comment{loop for input-enabledness}
    \end{algorithmic}
\end{transys}

\begin{transys}
    \label{PLTS:SRSD}
    \caption{Encoding of \textsf{SRSD}}
    \begin{algorithmic}
        \State \textbf{State:} $Asg.call, Asg.ret, Rec.call, Rec.ret \in \bool^{[n]};
        val, guess \in X \cup \{\bot\}; F \subseteq [n]$
        \State \textbf{Label:} $\{\text{Asg.call}(id), \text{Asg.return}(id), \text{Rec.call}(id),
        \text{Rec.return}(id,x), \text{fail}(id), guess(x) \mid id \in [n], x \in X\}$
        \State \textbf{Input:} $\{\text{Asg.call}(id), \text{Rec.call}(id), \text{fail}(id) \mid
        id \in [n]\}$
        \State \textbf{Fairness:} all transitions are fair except those labelled by fail and the
        call loops
        \State \textbf{Initial:} $Asg.call, Asg.ret, Rec.call, Rec.ret = \bot^{[n]};
        val, guess = \bot; F = \emptyset$
        \medskip
        \State
        $Asg.call[id] = \bot
        \xrightarrow{\text{Asg.call}(id)}
        Asg.call[id] = \top$
        \State
        $\top
        \xrightarrow{\text{Asg.call}(id)}
        \top$ \Comment{loop for input-enabledness}
        \State
        $|\{id \in [n] \setminus F \mid call[id] \ne \bot\} \cup F| \ge n-f
        \wedge Asg.ret[id] = \bot
        \xrightarrow{\text{Asg.return}(id)}
        Asg.ret[id] = \top$
        \State
        $Rec.call[id] = \bot
        \xrightarrow{\text{Rec.call}(id)}
        Rec.call[id] = \top$
        \State
        $\top
        \xrightarrow{\text{Rec.call}(id)}
        \top$ \Comment{loop for input-enabledness}
        \State
        $|\{id \in [n] \setminus F \mid call[id] \ne \bot\} \cup F| \ge n-f
        \wedge val = \bot
        \xrightarrow{\tau}
        val = x \text{ with probability } |X|^{-1}$
        \State
        $val \ne \bot \wedge Rec.ret[id] = \bot
        \xrightarrow{Rec.return(id,val)}
        Rec.ret[id] = \top$
        \State
        $id \not\in F \wedge |F| < f
        \xrightarrow{\text{fail}(id)}
        F = F \cup \{id\}$
        \State
        $\top
        \xrightarrow{\text{fail}(id)}
        \top$ \Comment{loop for input-enabledness}
        \State
        $\forall id \in [n] \setminus F, Rec.call[id] = \bot \wedge guess = \bot
        \xrightarrow{\text{guess}(x)}
        guess = x$
    \end{algorithmic}
\end{transys}

\begin{transys}
    \label{PLTS:BRB}
    \caption{Encoding of \textsf{BRB}}
    \begin{algorithmic}
        \State \textbf{State:} $call \in M\cup\bool; ret \in \bool^{[n]}; F \subseteq [n]$
        \State \textbf{Label:} $\{\text{call}(m), \text{return}(id,x), \text{fail}(id) \mid
        id \in [n], m \in M\}$
        \State \textbf{Input:} $\{\text{call}(m), \text{fail}(id) \mid id \in [n], m \in M\}$
        \State \textbf{Fairness:} all transitions are fair except those labelled by fail and the
        call loops
        \State \textbf{Initial:} $call = \bot; ret = \bot^{[n]}; F = \emptyset$
        \medskip
        \State
        $call = \bot
        \xrightarrow{\text{call}(m)}
        call = m$
        \State
        $\top
        \xrightarrow{\text{call}(id,b)}
        \top$ \Comment{loop for input-enabledness}
        \State
        $\text{leader is in } F
        \wedge ret = \bot^{[n]}
        \xrightarrow{\tau}
        call = \top$
        \State
        $call \in M \wedge ret[id] = \bot
        \xrightarrow{\text{return}(id,call)}
        ret[id] = \top$
        \State
        $id \not\in F \wedge |F| < f
        \xrightarrow{\text{fail}(id)}
        F = F \cup \{id\}$
        \State
        $\top
        \xrightarrow{\text{fail}(id)}
        \top$ \Comment{loop for input-enabledness}
    \end{algorithmic}
\end{transys}

\begin{transys}
    \label{PLTS:AVSS}
    \caption{Encoding of \textsf{AVSS}}
    \begin{algorithmic}
        \State \textbf{State:} $Sh.call, Rec.ret \in X \cup \bool; Sh.ret, Rec.call \in \bool^{[n]};
        guess \in X \cup \{\bot\}; F \subseteq [n]$
        \State \textbf{Label:} $\{\text{Sh.call}(x), \text{Sh.return}(id,x), \text{Rec.call}(id),
        \text{Rec.return}(id,x), \text{fail}(id), guess(x) \mid id \in [n], x \in X\}$
        \State \textbf{Input:} $\{\text{Sh.call}(id), \text{Rec.call}(id), \text{fail}(id) \mid
        id \in [n]\}$
        \State \textbf{Fairness:} all transitions are fair except those labelled by fail and the
        call loops
        \State \textbf{Initial:} $Sh.call = \bot; Sh.ret, Rec.call, Rec.ret = \bot^{[n]};
        val, guess = \bot; F = \emptyset$
        \State \textbf{Partial Information:} $\ocal_Q$ erases the valuation of $Sh.call$ if it is
        in $X$ and $\ocal_L$ removes $x$ in $Sh.call(x)$
        \medskip
        \State
        $Sh.call = \bot
        \xrightarrow{\text{Sh.call}(x)}
        Sh.call = x$
        \State
        $\top
        \xrightarrow{\text{Sh.call}(x)}
        \top$ \Comment{loop for input-enabledness}
        \State
        $\text{dealer is in } F
        \wedge Sh.ret = \bot^{[n]}
        \xrightarrow{\tau}
        Sh.call = \top$ \Comment{Bad dealer prevents termination of \textsf{Sh}}
        \State
        $\text{dealer is in } F
        \wedge Sh.ret = \bot^{[n]}
        \xrightarrow{\tau}
        Sh.call = x$ \Comment{Bad dealer commits an arbitrary value}
        \State
        $Sh.call \in X \wedge ret[id] = \bot
        \xrightarrow{Sh.return(id,Sh.call)}
        ret[id] = \top$
        \State
        $Rec.call[id] = \bot
        \xrightarrow{\text{Rec.call}(id)}
        Rec.call[id] = \top$
        \State
        $\top
        \xrightarrow{\text{Rec.call}(id)}
        \top$ \Comment{loop for input-enabledness}
        \State
        $|\{id \in [n] \setminus F \mid Rec.call[id] \ne \bot\} \cup F| \ge n-f
        \wedge Rec.ret[id] = \bot
        \xrightarrow{\tau}
        Rec.ret[id] = \top \text{ with probability }\epsilon; x \in X
        \text{ with probability } (1-\epsilon)*\epsilon; Sh.call \text{ otherwise}$
        \State
        $Rec.ret[id] \in X
        \xrightarrow{Rec.return(id,Rec.ret[id])}
        Rec.ret[id] = \top$
        \State
        $id \not\in F \wedge |F| < f
        \xrightarrow{\text{fail}(id)}
        F = F \cup \{id\}$
        \State
        $\top
        \xrightarrow{\text{fail}(id)}
        \top$ \Comment{loop for input-enabledness}
        \State
        $\forall id \in [n] \setminus F, Rec.call[id] = \bot \wedge guess = \bot
        \xrightarrow{\text{guess}(x)}
        guess = x$
    \end{algorithmic}
\end{transys}

\clearpage
\section{Proof of Theorem \ref{thm:Segala}}
\label{app:ProofSegala}

In our setting, the most interesting implication of Theorem \ref{thm:Segala} is soundness.
We will see that the definition of $\sqsubseteq^t$ has some limitations, notably with respect to
liveness. To address this, we introduce a stronger notion of trace distribution and strengthen
probabilistic forward simulation accordingly. However, the soundness proof of Segala
\cite{Segala1995Compositional} works at the level of $\efrak_{\mu,s}$, not $s$. This is
problematic because our strengthening consists in constraining schedulers; thus, we need to
adapt the proof. If $\scal$ is a probabilistic forward
simulation then it induces a choice function
$\ccal : \scal \times T \rightarrow Sch \times (\scal \rightarrow [0,1])$
which maps any $(q,\mu') \in \scal$ and $(q,l,\mu) \in T$ to the scheduler that induces
$\mu' \xRightarrow{l} \mu''$
and the witness $w$ of $(\mu,\mu'') \in \dcal(\scal)$.

\begin{algorithm}
    \caption{Simulation Algorithm for Schedulers}
    \label{alg:simulation}
    \begin{algorithmic}
        \State \textbf{Input:} $\pcal, \pcal'$ two PLTSs, $\scal$ a probabilistic forward
        simulation of $\pcal$ by $\pcal'$ and $s \in Sch$
        \medskip
        \State $\texttt{e}, \texttt{e'}, \texttt{mu'} \gets q_\star, q'_\star, \delta_{q_\star'}$
        \State $\texttt{acc} \gets \texttt{e'}$
        \State $\texttt{t} \gets s(\texttt{e})$
        \While{$\texttt{t} \ne \bot$}
            \State $\texttt{s'}, \texttt{w} \gets \ccal(last(e),\mu',t)$
            \State $\texttt{t'} \gets \texttt{s'}(\texttt{e'})$
            \While{$\texttt{t'} \ne \bot$}
                \State $\texttt{e'} \gets \texttt{e'}.lbl(\texttt{t'}).out(\texttt{t'})$
                \State $\texttt{acc} \gets \texttt{acc}.lbl(\texttt{t'}).out(\texttt{t'})$
                \State $\texttt{t'} \gets \texttt{s'}(\texttt{e'})$
            \EndWhile
            \State $\texttt{e'}, \texttt{mu'} \gets last(\texttt{e'}),
            \qfrak_{\texttt{mu'},\texttt{s'}}$
            \State $\texttt{e} \gets \texttt{e}.lbl(\texttt{t}).q$ with probability
            $\frac{\sum_{\mu \in \dcal(Q')}w(q,\mu)\cdot\mu(\texttt{e'})}
            {\texttt{mu'}(\texttt{e'})}$
            \State $\texttt{t} \gets \texttt{s}(\texttt{e})$
        \EndWhile
    \end{algorithmic}
\end{algorithm}

Once the choice function $\ccal$ is fixed, Algorithm \ref{alg:simulation} is a sequential
probabilistic program; thus, it can be modelled as a deterministic PLTS. Of course, the
validity of the algorithm's individual instructions should be established rigorously, but
this is not our concern here. The variable $\texttt{e}$ encodes an execution of $\pcal$ and
$\texttt{acc}$ an execution of $\pcal'$ that simulates it. In order to extract a scheduler of
$\pcal'$ from Algorithm \ref{alg:simulation}, we need to identify the instructions deciding which
transitions to execute in $\pcal'$. The most obvious one is
$\texttt{t'} \gets \texttt{s'}(\texttt{e'})$
when $\texttt{t'} \ne \bot$; then there is $\texttt{t} \gets \texttt{s}(\texttt{e})$ when
$\texttt{t} = \bot$. The case where $\texttt{t'} = \bot$ has to be discarded
because another transition will be picked without updating $\texttt{acc}$. When $\texttt{t} = \bot$,
the algorithm stops, so the scheduler on $\pcal'$ also returns $\bot$. After each update
of $\texttt{acc}$, one of the three following events occurs: (i) $\texttt{t}$ is set to $\bot$;
(ii) $\texttt{t'}$ is set to a non-$\bot$ value; or (iii) $\texttt{t'}$ is set to $\bot$
indefinitely.
Cases (i) and (iii) correspond to scheduling $\bot$ and case (ii) corresponds to scheduling
$\texttt{t'}$. Thus, the key values to compute are $reach(e')$, the probability of reaching a
configuration where $\texttt{acc} = e'$, and $sch(e',t)$, the probability of reaching a
configuration where $\texttt{acc} = e$ and $\texttt{t'} = t$. Knowing these, one can define a
scheduler on $\pcal'$ as
\[
    s' : e' \mapsto
        \begin{cases}
            t' \mapsto \frac{sch(e',t')}{reach(e')} \\
            \bot \mapsto 1 - \frac{\sum_{t' \in T'}sch(e',t')}{reach(e')}
        \end{cases}
\]
We conjecture that $\tfrak_{q_\star,s} = \tfrak_{q_\star',s'}$.
